\subsection{Gleichmächtigkeit eines Funktionsgraphen}

\Needspace{16\baselineskip}
\FormulaThmDeltaK[Ein Funktionsgraph ist zum Definitionsbereich gleichmächtig]{%
  F\colon D\to A\vdash D\EqCard F
}{FunctionGraphEquinumerousDomain}{%
  \DeltaRow{Mengen}{A\dsep D\dsep F\dsep i\dsep j\dsep a\dsep p}
  \DeltaRow{Abkürzung}{e=\{(i,p)\in D\times F\mid p=(i,F(i))\}}
}
Die Hilfsfunktion ordnet jedem Argument das zugehörige Paar des Graphen zu.
Ihre Injektivität beruht auf der ersten Komponente, nicht darauf, dass
verschiedene Argumente verschiedene Funktionswerte haben.
\begin{tabproofwide}
  \proofstepwidestar[1]{F\colon D\to A}{\rA}
  \proofstepwidestar[2]{i\in D}{\rA}
  \proofstepwidestar[1,2]{(i,F(i))\in F}{%
    \FormulaRefAuto{F\colon A\to B\dsep x\in A\vdash(x,F(x))\in F}[thm]{1,2}}
  \proofstepwidestar[1]{\forall i\in D\,((i,F(i))\in F)}{\rUI{\rRI{2,3}}}
  \proofstepwidestar[1]{\begin{aligned}[t]
    &e\colon D\to F\land{}\\[-2pt]
    &\forall i\in D\,(e(i)=(i,F(i)))
  \end{aligned}}{%
    \rAI{\FormulaRefAuto{TypedTermGraphFunction}[thm]{4},
      \FormulaRefAuto{TypedTermGraphValues}[thm]{4}}}
  \proofstepwidestar[6]{i\in D}{\rA}
  \proofstepwidestar[7]{j\in D}{\rA}
  \proofstepwidestar[8]{e(i)=e(j)}{\rA}
  \proofstepwidestar[1,6]{e(i)=(i,F(i))}{\rRE{6,\rUE{\rAEb{5}}}}
  \proofstepwidestar[1,7]{e(j)=(j,F(j))}{\rRE{7,\rUE{\rAEb{5}}}}
  \proofstepwidestar[1,6,7,8]{(i,F(i))=(j,F(j))}{\rIE{9,10,8}}
  \proofstepwidestar[1,6,7,8]{i=j}{\rAEa{\FormulaRefAuto{(a,b)=(c,d)\eqvdash a=c\land b=d}[thm]{11}}}
  \proofstepwidestar[1,6,7]{e(i)=e(j)\rightarrow i=j}{\rRI{8,12}}
  \proofstepwidestar[1,6]{\forall j\in D\,(e(i)=e(j)\rightarrow i=j)}{\rUI{\rRI{7,13}}}
  \proofstepwidestar[1]{\forall i,j\in D\,(e(i)=e(j)\rightarrow i=j)}{\rUI{\rRI{6,14}}}
  \proofstepwidestar[1]{e\colon D\inj F}{\FormulaRefAuto{Injektive Funktion}[def]{\rAEa{5},15}}
  \proofstepwidestar[17]{p\in F}{\rA}
  \proofstepwidestar[1]{F\subseteq D\times A}{\FormulaRefAuto{FunctionGraphSubsetProduct}[thm]{1}}
  \proofstepwidestar[1,17]{p\in D\times A}{%
    \FormulaRefAuto{A\subseteq B\dsep x\in A\vdash x\in B}[thm]{18,17}}
  \proofstepwidestar[1,17]{\exists i\in D\,\exists a\in A\,p=(i,a)}{%
    \FormulaRefAuto{z\in A\times B\eqvdash\exists a\in A\exists b\in B\;z=(a,b)}[thm]{19}}
  \proofstepwidestar[21]{i\in D\land\exists a\in A\,p=(i,a)}{\rA}
  \proofstepwidestar[22]{a\in A\land p=(i,a)}{\rA}
  \proofstepwidestar[17,22]{(i,a)\in F}{\rIE{\rAEb{22},17}}
  \proofstepwidestar[1,17,22]{F(i)=a}{%
    \FormulaRefAuto{F\colon A\to B\dsep(x,y)\in F\vdash F(x)=y}[thm]{1,23}}
  \proofstepwidestar[1,21]{e(i)=(i,F(i))}{\rRE{\rAEa{21},\rUE{\rAEb{5}}}}
    \proofstepwidestar[1,17,21,22]{e(i)=p}{\rIE{24,\FormulaRefAuto{x=y\vdash y=x}[thm]{\rAEb{22}},25}}
  \proofstepwidestar[1,17,21,22]{\exists i\in D\,e(i)=p}{\rEI{\rAI{\rAEa{21},26}}}
  \proofstepwidestar[1,17,21]{\exists i\in D\,e(i)=p}{\rEE{\rAEb{21},22,27}}
  \proofstepwidestar[1,17]{\exists i\in D\,e(i)=p}{\rEE{20,21,28}}
  \proofstepwidestar[1]{\forall p\in F\,\exists i\in D\,e(i)=p}{\rUI{\rRI{17,29}}}
  \proofstepwidestar[1]{e\colon D\sur F}{%
    \FormulaRefAuto{Surjektive Funktion}[def]{\rAEa{5},30}}
  \proofstepwidestar[1]{e\colon D\bij F}{%
    \FormulaRefAuto{Bijektive Funktion}[def]{16,31}}
  \proofstepwidestar[1]{D\EqCard F}{%
    \FormulaRefAuto{A\EqCard B\coloneqq\exists F\colon A\bij B}[def]{\rEI{32}}}
\end{tabproofwide}
